\section{Discussion}
\label{sec:ch3-discussion}

As an extension of HCI~\cite{alavi2019introduction}, HBI grounds our contribution of an AfI-driven approach to examining experiences in built environments. Below, we consider implications of our findings for smart building design, outline methodological extensions for HCI/HBI, reflect on AfI contributions, identify key tensions, and note limitations.

\subsection{From Experience to Design: Affective Interactions in Smart Buildings}
Our findings complement and contrast with current approaches in smart buildings, pointing to opportunities for systems that better align with how people actually engage with these environments.

Smart buildings are typically designed to detect and respond once a predefined threshold is breached (e.g. temperature exceeding 25°C)~\cite{jazizadeh2014human, marques2020real, de2002thermal}. In contrast, our participants adjusted their behaviour well before such thresholds were met, using bodily signals like muscle tension, changes in peripheral light, or feelings of exposure (W1, W2; AP1; IQ1). These anticipatory responses never triggered system detection. This reveals a mismatch: current systems are reactive and event-driven, while people’s real-time engagement is early, subtle, and embodied. Addressing this calls for design objectives that support pre-emptive adjustment (IQ6).

Automation is often for reducing cognitive load through invisible adaptation~\cite{becerik2022ten}. Yet participants did not avoid cognitive engagement; they initiated it. Through gestures like waving at motion sensors, they engaged in informal experimentation to interpret smartness (W5; AP5; IQ5). This labour is currently rarely supported. We suggest adopting to partial legibility~\cite{lim2009and}, inviting users to probe, test, and build evolving understandings of system behaviour.

Building design also assumes isolated occupants, offering personal controls. However, our findings show that people continuously coordinated with others, modifying behaviour to avoid disrupting shared use, interpreting body language, or yielding to implicit social rules (W3; AP3; IQ3). This points to the need for interaction strategies that respect and surface social context, such as reversible, non-intrusive adjustments or ambient indicators of shared states~\cite{salim2012probing, yu2018delight}.

Finally, systems intervene when discomfort is detected, but participants described moments of alignment, where everything felt right and no interaction was needed (W5; AP6; IQ6). These ``settling in'' states were fragile yet valuable. Current systems often ignore them, treating non-interaction as inactivity. Such states need to be recognised as positive, requiring designs that support pacing, holding, and experiential continuity, not just intervention~\cite{alavi2017comfort}.


\subsection{Methodological Reflections for the HCI---HBI Research Community}
Methodologically, our work responds to two smart building needs: developing methods that can understand AI-driven building interactivity, and reconciling HCI’s methodological traditions with architectural complexity as research shifts from ``artefacts to environments’’~\cite{alavi2019introduction}. To address these needs, we introduce a \emph{``methodological package''}. \textit{Situated elicitation} prompts occupants to articulate their experience \emph{in place}, capturing moment-to-moment appraisals that arise from immediate spatial, social, and sensory conditions. Next, \textit{building walk-throughs} operationalise experience as something that unfolds over time: participants move through the spaces, pausing to comment, interpret, or demonstrate bodily responses. This reveals micro-adjustments, anticipatory behaviours, and interpretive negotiations with automation that static sensing or point-based surveys cannot surface. We also adopted environmental \textit{body mapping}~\cite{turmo2023towards} as an reflective activity within walk-throughs, enabling participants to consider and externalise sensations, tensions, and shifts. Providing multiple expressive formats (spoken, written, and drawn/coloured) supported inclusivity in how participants articulated affective experience. While not analysed as a primary dataset, body mapping functioned as a generative prompt that enriched participants’ accounts of lived affect.

Together, these techniques form a replicable and adaptable approach for examining interpretation, regulation, and sense-making in various technology-rich environments. They reinforce our broader contribution to HCI: extending experience-centred interaction research by operationalising AfI at environmental scale, and providing methodological means for studying affect as emergent and interactional. Affective Practices will vary across settings, but this stance provides a way to trace that variation. We see two routes through which other HCI domains can draw on this work: (1) adapting our \textbf{methods} to study affect in other environments, and (2) extending our \textbf{conceptual results} (APs and IQs) to inform design in adjacent HCI domains.

In \textbf{automotive HCI}, empathic cars~\cite{rao2023affective, ghosh2022exploring} aim to sense and respond to driver and passenger emotion through lighting, sound, and haptic modulation~\cite{liu2021empathetic, nadri2022empathic}. Our methodological approach could be adapted into ``affective drives’’, ‘‘pause-and-reflect checkpoints’’, and in-car body mapping to trace emotional changes across journeys. Conceptually, practices such as \textit{Orient and Appraise}, \textit{Modulate and (De)Escalate} and \textit{Settle and Stabilise}, supported by \textit{Smartness Legibility} and \textit{Active Non-Intervention}, offer vocabulary for designing systems that regulate emotion, manage transitions in automation, and know when not to intervene.

Similarly, \textbf{virtual reality (VR)} research orchestrates atmosphere, pacing, proximity, and social cues, paralleling building-scale interaction~\cite{albarrak2023using, ahmadpour2025affectivevr}. Methodologically, we see ``VR affect walks'' or in-VR sense-making pauses, where participants trace their somatic and social responses as environments shift. Just as building walks captured temporal alignment/misalignment, VR could use live annotation tools, movement-logging, or post-immersion reconstructions to surface how emotional rhythms unfold inside simulated architectures. Conceptually, \textit{Spatial Niche Discovery} and \textit{Social Attunement and Coordination} offer guidance for environments that scaffold orientation, co-presence, and emotional pacing rather than focusing solely on immersion or task success.




\subsection{Reflections for the Affective Interaction Research Community}
\label{sec:ch3-disc-reflections-afi}
While the AfI perspective has emphasised situated and bodily meaning-making, its application has centred on digital or tangible systems (e.g. ~\cite{breazeal2001affective, zhou2024tangible}). Our study extends AfI into the built environment, highlighting a different interactional terrain~\cite{rao2025lessons}. Smart buildings are infrastructural, shared, and persistent systems that shape lived experience through slow, cumulative trajectories. By adopting an AfI lens, we surface under-explored experience in such environments (c.f. ~\autoref{sec:ch3-study2} and ~\autoref{sec:ch3-higher-order}) and offer several reflections for the AfI community.

\textit{First}, our experience using an AfI lens invites a move beyond event-based models of affect. In smart buildings, affective experience often emerges through bodily sensing (W1), interpretive hesitation (W2), spatial purpose-making (W3), social modulation and escalation (W4), and meaning-making around ‘smartness’ itself (W5). These are lived patterns that unfold over time and resist fixed classification. This perspective aligns with recent critiques of affect measurement in HCI~\cite{lottridge2011affective, ahmadpour2025affective}. It suggests that affective experience in built environments may be better understood through patterns of interpretation, reflection, and regulation than through momentary expressions alone.

\textit{Second}, the interaction qualities we propose (e.g. \textit{Active Non-Intervention}, \textit{Somatic Legibility}) for interactive smart buildings offer one way to support these temporal and interpretive dynamics. Rather than reacting to sensed emotion states, these qualities describe how systems can modulate presence or scaffold interpretation, without requiring overt expression or behavioural triggers. 

\textit{Lastly}, based on our experience of applying an AfI lens methodologically (through open-ended surveys and building walks) and analytically (through thematic synthesis), we found that AfI serves as a valuable framework for uncovering complex, temporally layered experiences that standard evaluation tools may overlook. While our prompts and structure around “experience” were intentionally open and interpretive, participants had little trouble engaging with them. Many navigated these vague questions with rich, situated narratives to articulate how they relate to their environment. This highlights AfI’s potential not just as a design stance, but as a methodological guide---one that can scaffold inquiry even without tightly defined constructs. It complements more structured smart building evaluation tools by revealing interpretive and affective couplings that might otherwise remain invisible.

\subsection{Tensions in Designing for Affective Interaction in Smart Environments}
As AfI enters the domain of smart buildings, new design tensions emerge, shaped by the infrastructural, shared, and persistent nature of these environments~\cite{rao2025lessons, rao2025designing}. Designing for AfI here requires more than sensing or reacting; it demands interpretive sensitivity for the subtle adjustments occupants already perform. Below, we outline and explain some of these tensions.

First, interaction qualities like \textit{Smartness Legibility} and \textit{Somatic Legibility} call for systems to surface their logic or state in ways occupants can interpret. But excessive signalling risks overwhelming users or disrupting ongoing tasks~\cite{case2015calm}. In shared environments, this legibility must be optional and timed, revealing just enough~\cite{mccarthy2004technology}.

Second, environments aim to be helpfully responsive, adjusting light, sound, or temperature in real time. Yet, when such support becomes too assertive, it can override user intention or suppress exploratory interaction~\cite{gaver2003ambiguity}. 

Next, ``affective fit'' is deeply personal, but buildings are inherently social. One's comfort may conflict with another’s focus. Our findings show that occupants often regulate themselves socially through deference, attunement, or negotiation, but current systems rarely reflect this nuance~\cite{erickson2000social, tolmie2002unremarkable, van2009social}. Designing for affect in shared spaces requires sensitivity to parallel user experience trajectories.

Finally, while our work rethinks affect, it does not reject affective computing or comfort-oriented models. Physiological sensing~\cite{picard1999affective} and affect mapping~\cite{russell1980circumplex} remain valuable, especially for detecting latent discomfort and gathering quantifiable metrics. We see these as complementary layers where affective computing provides system sensitivity, while AfI interaction design scaffolds the occupant’s meaning-making~\cite{mccarthy2004technology, hook2018designing}. Together, they form a more holistic response to affect in smart environments.


\subsection{Limitations}
\label{subsec:ch3-limitations}
Our study explored affective experience through emotion-centred dialogues with a small group of participants in a single smart building characterised by shared-use spaces and flexibility, including shifting occupancy patterns that shape how people negotiate comfort and meaning on a day-to-day basis. We posit our arguments within this scope. Nonetheless, our focus on lived, affective experiences offers varying degrees of transferable insights for future research. Future work could extend this across similar public and semi-public indoor typologies such as other university buildings, offices, classrooms, or public libraries, where spatial purpose, occupancy patterns, and shared infrastructures create conditions comparable to those that we studied. Replicating the studies in these settings would allow us to examine how \textit{Affective Practices} manifest under different uses, privacy, and environmental control while remaining within comparable organisational logics. However, our findings may not extend to privately owned residential spaces, where long-term settlement, personalisation, and individual control fundamentally reshape people’s affective orientations and adaptive strategies. Practices such as \textit{Spatial Niche Discovery}, for instance, take on different meanings when users can modify and claim space over months or years. Recognising these differences is important: tracing these practises across different settings can help clarify which affective processes travel across contexts and which are deeply rooted in specific spatial or institutional conditions. We see this comparative analysis as an important direction for future.

While our methodology supported rich, open-ended reflections, it is limited in scalability. Building walks requires time and trust, limiting broader deployment. We also prioritised subjective meaning-making over long-term sensing. This offers only partial insight into how affective regulation unfolds over extended periods. Future work could incorporate physiological or behavioural data to capture broader rhythms of engagement and scale using hybrid approaches such as mobile probes or spatial diaries. Finally, our \textit{Affective Practices} and \textit{Interaction Qualities} are exploratory. Further work is needed to understand their value in participatory design, inclusive sensing, and evaluation. We see these as an invitation for future work to operationalise and expand these concepts across diverse sites and populations.
